% !TEX root = ../main.tex
\chapter{Theoretical framework and empirical strategy}

This chapter outlines the theoretical framework used for our analysis. We begin by adapting a heterogeneous-firm trade model to the context of individual remittance behavior, enstablishing the theoretical foundations of how a migrant's transfer decisions can be influenced by economic shocks.

\section{A Remittance Model with Heterogeneous Migrants}

To analyze the impact of shocks on remittance flows, we adapt the logic of heterogeneous-firm trade models to the context of remittances. Specifically, we consider a model where migrants differ in income and altruism intensity, and where remitting requires paying both a fixed access cost and a variable cost proportional to the amount remitted. Formally, consider migrant $i$, who resides in U.S. state $s$, originates from Mexican municipality $m$, and chooses remittances in period $t$. The migrant's utility function is given by:
\begin{equation}
    U_{imst}
    =
    \left[
    \alpha
    \left(c^L_{imst}\right)^{\frac{\sigma-1}{\sigma}}
    +
    (1-\alpha)
    \left(
    \phi_{im}
    \left[
    R_{imst}
    +
    {y}^O_{mt}
    \right]
    \right)^{\frac{\sigma-1}{\sigma}}
    \right]^{\frac{\sigma}{\sigma-1}}
\end{equation}
Here $\phi_{im} > 0$ denotes the altruism or obligation intensity parameter, analogous to the productivity parameter in \textcite{melitz2003}. This strictly positive parameter ensures that sending remittances can increase the migrant's utility by raising origin-household resources. The variable $c^L_{imst} \geq 0$ represents the migrant's consumption in the destination country, and $R_{imst} \geq 0$ is the amount sent in remittances by migrant $i$. The weight of local consumption is given by $\alpha \in (0,1)$, and the elasticity of substitution between local consumption and origin-household resources is given by $\sigma$. Origin-household total resources are therefore defined as $R_{imst} + y^O_{imst}$, where the presence of $y^O_{imst}$ introduces a role for origin-country economic conditions: lower perceived origin resources increase the marginal utility of remittances, generating an insurance motive consistent with \textcite{yang_choi_2007}.

\par 

The migrant's budget constraint is given by:
\begin{equation}
    c^L_{imst}
    +
    f_{mt}\mathbbm{1}\{R_{imst} > 0\}
    +
    (1 + \tau_{mt})R_{imst}
    \leq
    w_{imst}
\end{equation}
where $w_{imst} \geq 0$ is the migrant's wage, $f_{mt} \geq 0$ is a fixed remitting cost capturing the pre-transaction expense of accessing a remittance channel, analogous to the fixed export cost in \textcite{melitz2003}, and $\tau_{mt} > -1$ is the variable cost which scales with the amount remitted, analogous to the iceberg cost structure in \textcite{melitz2003}. Notice that, even though $\tau_{mt}$ is typically positive to reflect transaction fees or exchange rate spreads, we allow for the possibility of negative variable costs, which could arise from the presence of remittance subsidies or favorable exchange rate regimes.

\par

If the migrant maximizes utility subject to the budget constraint and chooses to remit a positive amount ($R_{imst} > 0$), solving the Lagrangian yields the optimal remittance function:
\begin{equation}
R_{imst}^* = \frac{(w_{imst} - f_{mt}) - \frac{\phi_{im}}{A_{imst}} y^O_{mt}}{\frac{\phi_{im}}{A_{imst}} + (1 + \tau_{mt})}
\end{equation}
where $A_{imst}$ represents the optimal ratio of origin-household resources to own local consumption, and is defined as:
\begin{equation}
    A_{imst}
    \equiv
    \left(
    \frac{(1-\alpha)\phi_{im}}
    {\alpha(1+\tau_{mt})}
    \right)^{\sigma}
\end{equation}
To understand the effect of destination-country income shocks on remittance flows, we take the partial derivative of optimal remittances with respect to wage:
\begin{equation}
    \frac{\partial R^*_{imst}}{\partial w_{imst}}
    =
    \frac{A_{imst}}{1 + (1 + \tau_{mt})A_{imst}}
    >
    0
\end{equation}
Since the derivative is strictly positive, through an income effect the model predicts that a positive shock to the migrant's wage will increase remittances, while a negative shock will decrease them. The fixed cost $f_{mt}$ establishes an entry threshold: if the migrant's wage is not sufficiently high relative to the fixed cost, they will not send remittances at all. Additionally, the variable cost $\tau_{mt}$ reduces the marginal benefit of each dollar sent, dampening the sensitivity of remittances to wage fluctuations.

\section{Allowing for Spillovers}

We expand the model to allow for spillover effects across migrant networks. Decomposing origin-household resources into local non-remittance income and remittances received from migants in other US states, the perceived origin-household resources excluding migrant $i$'s own remittances can be expressed as:
\begin{equation}
    y^O_{imst} = \tilde{y}^O_{mt} + \sum_{s' \neq s} R_{s'mt}
\end{equation}
Solving for optimal remittances in the augmented model yields:
\begin{equation}
R_{imst}^* = \frac{(w_{imst} - f_{mt}) - \frac{\phi_{im}}{A_{imst}} \left(\tilde{y}^O_{mt} + \sum_{s' \neq s} R_{s'mt} \right)}{\frac{\phi_{im}}{A_{imst}} + (1 + \tau_{mt})}
\end{equation}
Hence, holding the migrant's own wage and remitting costs fixed, remittances from other migrants enter the decision through perceived origin-household resources:
\begin{equation}
\frac{\partial R_{imst}^*}{\partial \left( \sum_{s' \neq s} R_{s'mt} \right)} = - \frac{\phi_{im}}{\phi_{im} + A_{imst}(1 + \tau_{mt})}
\end{equation}
Because the derivative is strictly negative, high exposure to a destination shock lowers perceived origin resources and raises the marginal utility of alternative transfers. This drives compensatory network spillovers, where the direct remittance drop generates an offsetting increase from other destinations.

\par
To ground these concepts, the model delivers two empirical implications that guide our analysis. First, because optimal remittances are increasing in migrant income, a negative wage or income shock in the destination country should reduce remittances sent by affected migrants. Second, because remittances from other destinations enter the migrant's decision through perceived origin-household resources, a decline in transfers from the affected destination raises the marginal value of remittances from migrants in alternative destinations connected to the same municipality. Hence, origin municipalities with greater pre-existing exposure to the affected migrant destination should experience larger declines in total remittance receipts, all else equal. Based on these predictions, we develop our empirical strategy as follows: first, we test whether local inflows fall more in Mexican municipalities highly exposed to the destination hit by the shock; then, we examine whether co-exposure to other high-sending US states not directly affected mitigate the negative impact on total remittance revenue.
